% !TeX root = ../papers/nc-p3.tex
\section{一点を指定した反転窓の標準形と一般切断階数}
\label{sec:marked-opposite-rank}
\begin{theorem}\label{thm:marked-opposite-eight}
定理\ref{thm:reversible-component}の44次元成分の一般点の第三階数は8である。
従って14次元の有理商の幾何学的一般点では、標準AS正則性、最小周期2、第三階数8が同時に成立する。
\end{theorem}

隣接切断階数6の反転テンソルを、非退化対称対合を持つ六次元空間$H$への双線型写像
$\beta:E^*\times F^*\to H$のGramテンソルで表す。二重交代化の消滅条件は
\[
 \langle\beta(e,f),\beta(e',f')\rangle=\langle\beta(e,f'),\beta(e',f)\rangle .
\]
零点対$([e_0],[f_0])$では$\beta(e_0,f_0)=0$である。
$U=\beta(e_0,F^*)$と$V=\beta(E^*,f_0)$が三次元で、$U$の対合が非退化だとする。
両者は直交し、$H=U\perp V$となる。直交標準基底と一次空間の補空間を取ると、表示は
\[
 \beta((r,x),(s,y))=(r y+A_x y,\ s x+B_y x),\qquad x\in V,\quad y\in U
\]
となる。一文字に$e_0$または$f_0$を入れたGram等式は各$A_x,B_y$の対称性を与える。
補空間を変えると各作用素に恒等写像の線型なスカラー倍を加えられ、跡を零にできる。

三つの三行対称行列$X=(X_1,X_2,X_3)$について、巡回順$(23),(31),(12)$を使い
$D_X$の第$(ij)$列、第$(mn)$行を$[X_i,X_j]_{mn}$とする。
上の$A,B$を基底上の三行列として読むと、残る九つのGram等式はすべて
\[
 D_B=-D_A^{\mathsf t}
\]
にまとめられる。これは第$(ik),(jl)$の式が
$[A_i,A_k]_{jl}+[B_j,B_l]_{ik}=0$であることから従う。

\begin{lemma}[三行の交換子と直交二次式]\label{lem:ternary-net-commutator}
対称行列$S$が$\operatorname{tr}S=\operatorname{tr}(SX_i)=0$を満たすとき、
$P_X(S)_{ij}=\operatorname{tr}(SX_iX_j)$とおけば
\[
 D_X P_X(S) D_X^{\mathsf t}=(\det D_X) S .
\]
\end{lemma}
\begin{proof}
$P_X$は転置と跡の巡回性から対称である。各$X_i$へのスカラー行列の加法は両辺を変えない。
三行列の並べ方を列変換$g$で変えると$D$は右から$(\det g)g^{-\mathsf t}$を掛けられ、
$P$は$g^{\mathsf t}Pg$となるので等式は共変である。三次元空間の向きを保つ直交変換にも共変である。
非零な無跡対称$S$を動かしたとき$\operatorname{tr}(SX_i)=0$の三つ組はベクトル束をなし、
この全空間は既約である。従って相異なる固有値を持つ一般の$S$と一般の三つ組で証明すれば足りる。
$S=\operatorname{diag}(p,q,r)$、$p+q+r=0$とし、
$L=\operatorname{diag}(q-r,r-p,p-q)$とする。基底変換とスカラーの加法により
$X_i$の非対角部分を順に$E_{23}+E_{32},E_{31}+E_{13},E_{12}+E_{21}$、
対角部分を順に$aL,bL,cL$とできる。
$v=(a,b,c)^{\mathsf t}$、$[v]_\times$を$v$との外積の行列とすると直接乗算で
\[
 D=-I+3S[v]_\times,\qquad
 P=-S+\operatorname{tr}(SL^2)vv^{\mathsf t}.
\]
この二式に$r=-p-q$を代入して$DPD^{\mathsf t}-(\det D)S$を展開すると全九成分は零である。
\path{computations/verify_marked_opposite_rank.py}は同じ計算を五つの不定元を残して検査する。
$S=0$での等式も直ちに従う。
\end{proof}

\begin{proof}[定理\ref{thm:marked-opposite-eight}の証明]
まず$D_A$が可逆な表示を考える。対称対の最初の四つを$00,01,02,03$とすれば、
反対切断の十行十列行列は$\left(\begin{smallmatrix}0&L_A\\L_B^{\mathsf t}&W\end{smallmatrix}\right)$である。
上の四行は六次元の対称行列空間で$I,A_1,A_2,A_3$、左の四列は$I,B_1,B_2,B_3$に対応する。
$D_A,D_B$が可逆ならこれらは各々独立である。両側の直交補の行列をそれぞれ$S,T$とすると、
消去後の残り二行二列の対合は
\[
 \operatorname{tr}(T P_A(S))+\operatorname{tr}(S P_B(T))
\]
である。実際、右下の$ik,jl$成分は$(A_iA_k)_{jl}+(B_jB_l)_{ik}$の対称対への制限である。
補題から$P_A(S)=\det(D_A)D_A^{-1} S D_A^{-\mathsf t}$となり、
$D_B=-D_A^{\mathsf t}$と列数3の符号で二つの跡が打ち消し合う。従って第三階数はちょうど8となる。

次にこの開条件が目的の成分の一般点で成立することを確かめる。
二十点例の第一の零点対を指定する。点を指定した有限エタール被覆の完備局所環は
元の点のものと同じである。$P,H,Y,G,q$は定理\ref{thm:reversible-component}のGram表示とする。
選んだ$e_0,f_0$を標準行の先頭へ置き、各々に残り三つの座標行を補う。
十六行の変換後の$P$のうち$(0,j)$の三行を$U_0$、$(i,0)$の三行を$V_0$と書く ($i,j>0$)。
$J=U_0 H U_0^{\mathsf t}$は可逆である。$(i,1),(i,2),(i,3)$の行列を$P_i$とすると、
直交化前には$C_i=P_i H U_0^{\mathsf t}$が対称で、交換子は
$C_iJ^{-1}C_j-C_jJ^{-1}C_i$で計算する。基底を直交化してもその三列の独立性は変わらない。

例ではこの交換子行列は$\operatorname{diag}(-15/32,0,0)$である。
$\ker G$上の変分で点も動かし、可逆な最初の一行を消去した残り二行二列の微分を取る。
具体的には$\dot P=YH^{-1}$に対して六つの零点条件を微分し、
$V_0,U_0$を重ねた六行六列の転置で割って点の六座標の変分を得る。
上記ファイルの有理な接線の直線上では、接錐の方程式と二行二列の微分の行列式は順に
\[
 -(928t^2+704t+232)/15,\qquad (112t^2+184t-37)/256 .
\]
両式は互いに素なので、後者の斉次二次式は$q$の倍数でない。
完備局所環の随伴次数環が$(q)$で割った整域であることから、この小行列式の二次初項は消えず、
交換子行列の行列式は成分上で恒等的に零にならない。従って一般点で使った開条件が成立する。
さらに特殊ファイバーの二十の印は対角$S_5$と両因子交換で推移的に移り、
両因子交換は同じ成分と可逆条件を保つ。よってこの成分の一般点では二十通りの印のどれも使える。
\end{proof}

\begin{proposition}[印付きの一般窓を構成するパラメータ]\label{prop:marked-net-parameters}
上の可逆な開部分の印付きテンソルは、独立な無跡対称三行列$A$と、別の無跡対称三行列空間内の
三次元部分空間$\mathcal B$から構成できる。いずれの交換子行列も可逆なものを取る。
$\mathcal B$の基底$B^0$を取り、$K=D_{B^0}^{-1}(-D_A^{\mathsf t})$とすれば
\[
 B=B^0 g,\qquad g=\pm\sqrt{\det K}\,K^{-\mathsf t}.
\]
印を保つ同型ではこの二符号は同じである。正規化された二つの表示が同型となる条件は
$O_U\in\mathrm O(U), O_V\in\mathrm O(V), c\in k^\times$で
\[
 A'_{O_V x}=c O_U A_x O_U^{-1},\qquad
 B'_{O_U y}=c O_V B_y O_V^{-1}
\]
とできることである。従って上記14次元の有理商は単有理的である。
\end{proposition}
\begin{proof}
三列の基底変換の公式$D_B=D_{B^0}(\det g)g^{-\mathsf t}$を解くと表示式となる。
このとき九つのGram等式を満たし、定理の点付き開部分の全てのデータが得られる。
六次元空間への写像は全射であるから、Gramテンソルの射影同型は$H$での相似変換を誘導する。
指定の二直線を保つ基底変換を$(r,x)\mapsto(a r+u(x),C x)$、
$(s,y)\mapsto(b s+v(y),D y)$と書くと、相似変換の二つの直交因子への制限は
$aD=d O_U$, $bC=d O_V$と書ける。同型条件の跡を取ると、無跡性から$u=v=0$になる。
残る式は$c=ab/d$とした表示条件であり、逆にこの条件から基底変換を戻せる。
$O_U=-I,O_V=I,c=1$で$B$の符号だけを変えられる。
最初の関係空間に戻るときは二重交代化の核がただ一つの直線である開部分を使う。

$A$と$\mathcal B$のパラメータの基底はアフィン15空間と$\Gr(3,5)$の積の開部分で有理的である。
平方根を選ぶ被覆の二点は同じ軌道を与えるので、この基底から有理商への写像は有理関数体上に降下する。
目的の成分の幾何学的一般点はこの構成に含まれ、テンソルの局所成分は例で一つだから、
パラメータからの像の閉包はその成分の商となる（正則性を使うのは幾何学的一般点）。
これにより支配有理写像を得る。特殊な全データのAS性には別の判定が必要である。
\end{proof}
